\documentclass[../../main2.tex]{subfiles}

\begin{document}

\chapter{通道三：政治（经济 + 不可标价的偏好）}
\label{ch:politics}

\begin{motivation}
	上一章遗留的问题：经济通道的前提是「等价交换」——但你心里有几样东西\textbf{不许标价}：尊严、身份、公平。第 2 章说过，不可通约的冲突无解于算盘。那这些冲突谁来裁？\\
	\textbf{核心动机}：加入「不可货币化的偏好」这个约束，从经济里把政治推导出来——政治不是经济的上层建筑，是\textbf{交换失灵时接手裁决}的东西。\\
	\textbf{扩增模式}：有压，且加上身份权重——零和之外，还多了「谁配」的争夺。
\end{motivation}

\begin{logicmap}
\begin{tikzpicture}[node distance=0.85cm, every node/.style={draw, rectangle, rounded corners, align=center, font=\small}]
\node (q) {有些东西不许标价，谁来裁？};
\node (s) [below=of q] {不可货币化的偏好};
\node (sd) [below=of s] {认同赋予权力};
\node (h) [below=of sd] {现场：王安石与职场};
\node (c) [below=of h] {政治通道的账};
\node (n) [below=of c] {裁决怎么固化？（抛给第 14 章）};
\draw[-{Stealth}] (q) -- (s);
\draw[-{Stealth}] (s) -- (sd);
\draw[-{Stealth}] (sd) -- (h);
\draw[-{Stealth}] (h) -- (c);
\draw[-{Stealth}] (c) -- (n);
\end{tikzpicture}
\end{logicmap}

\section{交换里有些东西不卖}

\begin{readingnote}
	你有没有想过，为什么「给你加薪，这口锅你背了」这种话，很多人宁可辞职也不接？\\
	钱明明够多——是什么东西，钱买不动？
\end{readingnote}

第 2 章埋的线在这里炸响：有些偏好没有汇率。三个当场试验：

\begin{itemize}
	\item \textbf{尊严}：当众羞辱同事，赔多少钱「扯平」？——大多数人不接受任何数字。
	\item \textbf{身份}：「你说你是哪里人」的争论，能让两个高薪白领动手。GDP 解释不了它。
	\item \textbf{公平}：最后通牒实验——分一百块，你拿九十五我拿五，我宁可两块都不要，也要否掉你。\textbf{人会花钱买公平}。
\end{itemize}

用本书的标准句式：如果一切偏好都能标价，交换就能摆平一切冲突——世界只需要会计。但现在加入约束：\textbf{存在不可货币化的偏好}。这一步改变了所有性质：算盘失效的地方，必须有一个\textbf{不按报价说话}的裁决装置。这就是政治从经济里分离出来的时刻。

\begin{tcolorbox}[colback=green!5, title=\textit{生活类比}]
	经济像记账，政治像分家。账可以慢慢算，分家那天总有几样东西「不能折钱」——祖宅的堂屋、老人的赡养名分、谁在族谱头一页。算盘停转的地方，族长上场。族长的权力，来自「大家都认他能定这个」。
\end{tcolorbox}

\paragraph*{逻辑衔接}
	不可标价的偏好摆出来了。可裁决的凭什么听他的？下一节：权力的来源——认同。

\section{认同赋予权力}

\begin{readingnote}
	你有没有想过，「领导」两个字没有任何物理力量——为什么大家听他的？\\
	权力不是东西，是\textbf{一个关于「谁说了算」的共识}。
\end{readingnote}

权力的定义，一句话：\textbf{被认同赋权的裁决资格}。拆开是三步：

\begin{enumerate}
	\item \textbf{冲突不可标价}（上一节）——必须有人定；
	\item \textbf{众人默认某人能定}——认同；
	\item \textbf{定案被服从}——裁决落地。
\end{enumerate}

注意第二步：认同是可撤销的。「皇帝轮流做」的时刻，就是认同这个模因批量撤回的时刻——政治通道和模因通道在这里\textbf{咬合}：造反先造谣，夺权先夺话语权。

\begin{questionbox}
	\textit{现在你可能想反驳：「权力明明来自暴力和钱，讲认同太天真了。」}\\
	\textbf{暴力和钱是权力的燃料，不是引擎。}枪口要人抬，钱要人认。纯靠暴力的政权维持成本极高，随燃料断供垮台——历史上反复出现。「认同」不等于「心服」，只等于「按你说的办」——恐惧也算一种暂时的认同。我的用法是操作性的，不带温情。
\end{questionbox}

\paragraph*{逻辑衔接}
	权力的机制清楚了：裁不可标价之事，靠认同上台。但机制一进真实历史就变形——看两个现场：一场九百年前的改革，和一间今天的会议室。

\section{两个现场：王安石变法，和你楼下的晋升会}

\textbf{现场一：王安石变法（1069–1085）}。动机无可挑剔：富国强兵，民不加赋而国用饶。青苗法本意是官府低息贷给农民，压掉民间高利贷。为什么黄了？拆三层：

\begin{itemize}
	\item \textbf{渠道被劫持}：执行的还是那批士大夫与胥吏，贷款摊派成了政绩 KPI——低息贷变成了强制贷。动机好，渠道坏。
	\item \textbf{模因战打输}：「与民争利」的压缩版先入为主，司马光们的反对传得比新法的辩护快（第 11 章：强度压倒正确性）。
	\item \textbf{利益集团反噬}：高利贷的既得者、被摊派扰动的豪强，全成了反对派的兵源。政治通道里，\textbf{受损者的动员力}几乎总大于受益者的。
\end{itemize}

\textbf{现场二：职场晋升}。为什么「能力强」不等于「被提拔」？多因分账（第 9 章）走一遍：能力（A）、可见度（B）、站队（C）、部门预算（D）、老板心情（E）。做一次无你检验：拿掉 B——你干的活没人看见——晋升还会发生吗？多数岗位的答案是「不会」。结论刺耳但好用：\textbf{你优化 A 优化到 120 分，可能不如把 B 从 30 分提到 60 分}。学生看成绩排名，上班族看晋升，结构是一样的：\lowfreq{单一指标的排名，几乎永远是多因压缩后的失真投影。}

\begin{warningbox}
	别把这一节读成「躺平」或「厚黑」。它是\textbf{诊断}：想改变结果，先看哪一项的 $C$（贡献）最大、且你动得了。把全部筹码压在贡献 20\% 的变量上，才是真的不理性。
\end{warningbox}

\paragraph*{逻辑衔接}
	两个现场看完了：政治是不可标价冲突的裁决，靠认同运转，会被渠道和动员力改写。但裁决每次都吵一遍，社会受不了——下一章，法律出场：把裁决\textbf{写死}。

\section{政治通道的账}

这一通道的贡献怎么估？还是无你检验：「拿掉这场认同动员，结果会变吗？」三个快例：

\begin{itemize}
	\item \textbf{古}：拿掉「清君侧」的口号，七国之乱还打得起来吗？——会，但慢、散、弱。口号的 $C$ 为正。
	\item \textbf{今}：拿掉「涨价听证会」的民意压力，公用事业调价方案会改吗？——看 $C$ 大小，各地不同，可查可比。
	\item \textbf{今}：拿掉那次部门站队，晋升名单会变吗？——这就是你能在自己公司做的分账练习。
\end{itemize}

\begin{reader_anticipation}
	\item 权力靠「认同」，那认同本身怎么变？→ 认同是模因——第 11 章的传播规律管它，造反先造谣就是这么来的。
	\item 裁决每次都吵，社会怎么受得了？→ 把吵过的写死——第 14 章：法律作为定责编码登场。
\end{reader_anticipation}

\paragraph*{逻辑衔接}
	政治通道的账能记了。但每次裁决都重新吵一遍，成本太高——社会需要把「吵过的结果」固化下来。第 14 章：法律作为第四个齿轮，也是生成树的最后一环。

\begin{thinkbox}
\begin{enumerate}
	\item 回忆一次你在场的「不可标价」冲突（分家、评优、名额、背锅）：当时是钱说了算、认同说了算，还是力气说了算？用「裁决资格从哪来」重新看一遍。
	\item \textbf{反事实}：如果王安石的青苗法由一套不可被胥吏扭曲的系统执行（比如现代的直达补贴），变法会成吗？把「动机 / 渠道 / 模因」三笔账分个比例，再说说你对哪一笔最没把握。
	\item \textbf{反事实}：如果晋升完全公开打分（每一项权重、每次评分都透明）——不可标价的部分（信任、眼缘）被强行标价——办公室政治会消失，还是换一种玩法？
\end{enumerate}
\end{thinkbox}

\end{document}
